% Agent 16: Book pages 403--406
% Content: Section 4.7 (continuation) Gas Flow Near the Horizon,
%          Section 4.8 Accretion onto a Moving Hole (with Figure 4.8.1)

%%%%%%%%%%%%%%%%%%%%%%%%%%%%%%%%%%%%%%%%%%%%%%%%%%%%%%%%%%%%%%%%%%%%%%%%%%%%%%%
% Continuation of Section 4.7 -- Gas Flow Near the Horizon
%%%%%%%%%%%%%%%%%%%%%%%%%%%%%%%%%%%%%%%%%%%%%%%%%%%%%%%%%%%%%%%%%%%%%%%%%%%%%%%

% [Continuation from previous agent -- page 403 begins mid-section]

has the first integral
\begin{equation}
(dr/dt)^2 = 2GM/r \quad \text{(for fall from rest at } r \gg r_g\text{)}, \tag{4.7.2}
\end{equation}
and because the law of rest-mass conservation $\nabla \cdot (\rho \mathbf{u}) = 0$ has the first integral
\begin{equation}
4\pi r^2 \sqrt{-g}\, \rho\, (dr/dt) = \dot{M}_0, \tag{4.7.3}
\end{equation}
the density of rest mass must vary as
\begin{equation}
\rho = \frac{\dot{M}_0}{4\pi r^2} \left(\frac{r}{2GM}\right)^{1/2}
\simeq (6 \times 10^{-12} \;\text{g\,cm}^{-3})\, (r/r_g)^{-3/2}. \tag{4.7.4}
\end{equation}

This must be as true near and inside the horizon, $r \leq r_g$, as it is in the Newtonian
realm $r \gg r_g$. Similarly, the temperature will retain its Newtonian form (4.4.2):
\begin{equation}
T \simeq (10^{12}\;\text{K})\, (r_g/r); \tag{4.7.5}
\end{equation}
and the synchrotron emissivity will retain its Newtonian value (4.4.4). However,
the radiation which reaches infinity will be sharply cut off as the gas element
passes through the horizon.

\begin{equation}
\left(\text{fraction of emitted radiation that reaches } r = \infty
\text{ from infalling gas at radius } r_f\right)
\simeq \frac{27}{64}\left(1 - \frac{r_g}{r_f}\right). \tag{4.7.6}
\end{equation}

In previous sections we have taken rough account of this cutoff and of the
accompanying redshift by retaining the Newtonian luminosity down to $r = 2r_g$.

In the case of a rotating (Kerr) hole, near and inside the ergosphere the
dragging of inertial frames will swing the infalling gas into orbital rotation
about the hole. As the gas approaches the horizon, its angular velocity as seen
from infinity must approach the angular velocity of the horizon,
\begin{equation}
\Omega \to \Omega_{\text{horizon}} = \frac{a}{r_+^2 + a^2}
= \frac{c^3}{2GM} = \left(\text{sec}\right)^{-1}
\frac{10^5}{M_\odot}\left(\frac{M}{M_\odot}\right)^{-1}
\quad \text{for ``maximally rotating hole'', } a = M. \tag{4.7.7}
\end{equation}
(geometrized units, $c = G = 1$)

Here $a$ is the hole's angular momentum per unit mass. Any extra-luminous hot
spot in the gas (e.g.\ a region where magnetic field lines are reconnecting; cf.\
\S2.9) must orbit the hole with this angular velocity as it falls inward. The
brightness of the hole as seen at Earth will be modulated with a period
$P \simeq 4\pi r/\Omega \simeq (10^{-4}\;\text{sec})(M/M_\odot)$ as a result of the Doppler shift of the spot's
light and the bending of its rays. (A factor of $2\pi/\Omega$ is the orbital period; the
other factor of $2\pi/\Omega$ is the light travel time between the radius at which the
spot is located at the beginning of one period and the radius to which it has
fallen at the end of the period.) Thus, one might expect quasiperiodic
fluctuations in the light from a black hole living alone in the interstellar medium.

%%%%%%%%%%%%%%%%%%%%%%%%%%%%%%%%%%%%%%%%%%%%%%%%%%%%%%%%%%%%%%%%%%%%%%%%%%%%%%%
% Section 4.8 -- Accretion onto a Moving Hole
%%%%%%%%%%%%%%%%%%%%%%%%%%%%%%%%%%%%%%%%%%%%%%%%%%%%%%%%%%%%%%%%%%%%%%%%%%%%%%%

\subsection{Accretion onto a Moving Hole}
\label{sec:4.8}

Turn attention now from a hole at rest in the interstellar medium, to a hole that
moves with a speed much larger than the sound speed, $u_\infty \gg a_\infty$. Examine the
resulting accretion in the rest frame of the hole. At radii $r \gg 2GM/u_\infty^2$, the kinetic
energy per unit mass, $\frac{1}{2}u_\infty^2$, is far larger than the potential energy, $GM/r$; so the
pull of the hole has no effect on the gas flow. At $r \sim 2GM/u_\infty^2$, the pull of the
hole becomes significant. Because the flow is supersonic the gas will respond to
that pull in the same manner as would noninteracting particles; its pressure
cannot have an influence. (One can see this, for example, in the Euler equation
\begin{equation}
\text{const.} = \tfrac{1}{2}v^2 + \frac{a^2}{\Gamma - 1} - \frac{GM}{r}\,, \tag{4.8.1}
\end{equation}
where the speed-of-sound (enthalpy) term---which accounts for pressure effects---is
negligible compared to the kinetic-energy term.) Thus, the flow lines will be
identical to the trajectories of test particles: they will be hyperbolae
(Fig.~\ref{fig:4.8.1}a).

After passing around the hole, the trajectories of test particles intersect
(particle collisions; dotted line of Figure~\ref{fig:4.8.1}a). In the gas-dynamic case such
intersection of flow lines is prevented by rapidly mounting pressure. Con\-sequently,
a shock front must develop around the black hole (Fig.~\ref{fig:4.8.1}b).
Outside the shock front the flow lines will be hyperbolic test-particle trajectories.
Behind the shock front, they will not. The shock will be located roughly where
potential energy equals kinetic energy, so its characteristic size $l_s$ as shown in
Figure~\ref{fig:4.8.1}b will be
\begin{equation}
l_s \sim GM/u_\infty^2. \tag{4.8.2}
\end{equation}

In the shock, the gas will lose most of its velocity perpendicular to the shock
front [``strong shock'' in terminology of \S2.7, ``strong'' because $u_\infty/a_\infty = $
(speed of gas)/(speed of sound) $\gg 1$ on front side]; but it will retain all of its
velocity parallel to the front (denote this by $u_\|$). For a given gas element, if the
remaining kinetic energy behind the front greatly exceeds the potential energy,
$\frac{1}{2}u_\|^2 \gg GM/r$, then the gas element will escape the pull of the hole. If $\frac{1}{2}u_\|^2 \lesssim GM/r$,
then it cannot escape. The dividing line between escape and capture is a dotted
line in Figure~\ref{fig:4.8.1}b; and the corresponding impact parameter is labelled
$b_{\text{capture}}$.

We can calculate $b_{\text{capture}}$, in order of magnitude by idealizing the shock as
confined to the thin dotted region of Figure~\ref{fig:4.8.1}a, (Hoyle and Lyttleton 1939).
The value of $b_{\text{capture}}$ will correspond to an orbit for which
\begin{equation}
\tfrac{1}{2} v_\perp^2 = GM/x \quad \text{at} \quad y = 0. \tag{4.8.3}
\end{equation}

Straightforward examination of Kepler orbits reveals that
\begin{equation}
b_{\text{capture}} = 2GM/u_\infty^2. \tag{4.8.3'}
\end{equation}
Since the mass flux in the gas at large radii is $\rho_\infty u_\infty$, the accretion rate is
\begin{equation}
\dot{M}_0 = (\pi b_{\text{capture}}^2)\rho_\infty u_\infty
= 4\pi G^2 M^2 \rho_\infty / u_\infty^3. \tag{4.8.4}
\end{equation}

\begin{figure}[htbp]
\centering
% \includegraphics[width=\columnwidth]{fig_4_8_1}
\fbox{\parbox{0.9\columnwidth}{\centering [Figure 4.8.1 placeholder]\\
(a) Flow lines of supersonic gas past a black hole, treated as
test-particle trajectories. Dotted line shows region of particle
collisions (trajectory intersections).\\[6pt]
(b) Flow lines of supersonic gas (solid lines) showing the shock front
(heavy line). The characteristic size $l_s$ and capture impact parameter
$b_{\text{capture}}$ are indicated. The dividing line between escape and capture
is shown as a dotted line.}}
\caption{\textit{Figure 4.8.1.} The trajectories of test particles (a), and the flow lines of supersonic gas (b)
relative to a black hole, as viewed in the rest frame of the hole. The trajectories and flow
lines are virtually the same, until the gas encounters the shock front.}
\label{fig:4.8.1}
\end{figure}

More careful calculations give accretion rates which differ from this by multiplicative factors of $\sim 0.5$ to 1.0 (Bondi and Hoyle 1944).

Notice that the rate (4.8.4) for the supersonic case, and the rate
(4.2.10) for a hole at rest are identical except for the replacement of sound
speed $a_\infty$ by speed of flow $u_\infty$, and except for a numerical coefficient of order
unity. In the intermediate case of $u_\infty \sim a_\infty$, one can use the hybrid formula

\begin{equation}
\dot{M}_0 \simeq \frac{4\pi G^2 M^2 \rho_\infty}{(a_\infty^2 + u_\infty^2)^{3/2}}
\simeq \frac{(M/M_\odot)^2 \rho_\infty / (10^{-24}\;\text{g\,cm}^{-3})}
{(u_\infty / 10\;\text{km\,sec}^{-1})^2 + (a_\infty / 10\;\text{km\,sec}^{-1})^2]^{3/2}} \tag{4.8.5}
\end{equation}
(Bondi 1952).

[Salpeter (1964) gave the first application of these formulas to black holes; all
previous work dealt with normal stars.]

It is worth noting that a significant fraction of the kinetic energy released in
the shock
\begin{equation}
\frac{d\dot{E}}{dt} \simeq \tfrac{1}{2} u_\infty \dot{M}_0
\simeq \left(10^{24}\;\text{ergs}\right)
\left(\frac{M}{M_\odot}\right)^2
\left(\frac{\rho_\infty}{10^{-24}\;\text{g\,cm}^{-3}}\right)
\left(\frac{u_\infty}{10\;\text{km\,sec}^{-1}}\right)^{-1/2} \tag{4.8.6}
\end{equation}
may go into excitation of unionized atoms and thence into light as the atoms
decay, thus producing a shock front that glows (albeit weakly). See, e.g.,
Pikel'ner (1961), and \S\S\,10 of Kaplan (1966).

The motion of the gas behind the shock is practically unknown, particularly
when one takes account of magnetic fields. Many different types of instabilities
may occur, both here and in the case of a stationary black hole. For example at
radii $r < l_s$, where the inflow has become supersonic again, the Bisnovatyi-Kogan and
Sunyaev (1971) variety of plasma instabilities may also arise. And the reconnection
of magnetic field lines will surely be important. Almost all one can only guess that
the overall, time-averaged angular momentum per unit mass of the accreting
gas has no luminosity is produced, is angular momentum. Suppose that the accreting
Schwartzman (\S4.4).

One factor that may be important near the gravitational radius, where the most
of the gas luminosity is produced, is angular momentum in the direction of motion of the hole.
Such a solution is partially irrelevant because it ignores the intensity
($\sim$ direction in Fig.~\ref{fig:4.8.1}). If the angular momentum per unit mass, $L$, exceeds
$r_g c$, then centrifugal forces will become important before the infalling gas reaches
the horizon; the gas will be thrown into circulating orbits; and only after viscous
stresses have transported away the excess angular momentum will the gas fall
down the hole. (See \S5 for a situation similar to this.) The resulting viscous
heating and lengthened time spent outside the hole, as well as the changed flow
pattern, may affect the outpouring significantly. Moreover, when the
hole moves from a region with angular momentum of gas in one direction to a
region with angular momentum in another, the gas flow may become particularly
violent for a while near the horizon; and a strong flare of luminosity may be
produced (Salpeter 1964; Schwartzman 1971). Almost nothing is known today
about these issues.

Is the specific angular momentum of the accreting gas likely to exceed $r_g c$?
Interstellar gas is accreted from a region of diameter
\begin{equation}
2 b_{\text{capture}} = \frac{2r_g}{u_\infty^2} = (3 \times 10^{14}\;\text{cm})
\left(\frac{M}{M_\odot}\right)
\left(\frac{u_\infty}{10\;\text{km\,sec}^{-1}}\right)^{-2}. \tag{4.8.7}
\end{equation}
(For $u_\infty < a_\infty$, we must replace $u_\infty$ by $a_\infty$.) The interstellar gas is turbulent with
the turbulent velocity $v_{\text{turb}}$ on scales $l \sim b_{\text{capture}}$ given roughly by
\begin{equation}
v_{\text{turb}} \sim \left(\frac{\text{cm}}{\text{sec}}\right)
\left(\frac{10^{20}\;\text{cm}}{l}\right)^q
\left(3 \times 10^{20}\;\text{cm}\right)  \tag{4.8.8}
\end{equation}
where $q$ is a coefficient believed to lie between 1/2 and 1 (Kaplan and Pikel'ner
1970), and the coefficient $10^6$~cm/sec is taken from astronomical observations.
Let us take $q = 0.75$ as a best guess. Then the specific angular momentum
associated with this turbulence, for a gas element of size $l$, is
\begin{equation}
L_{\text{turb}} \simeq v_{\text{turb}} \, l
\simeq \left(3 \times 10^{26}\;\text{cm}^2\;\text{sec}^{-1}\right)
\left(\frac{l}{3 \times 10^{20}\;\text{cm}}\right)^{1.75}. \tag{4.8.9}
\end{equation}

Consequently, the specific angular momentum of the gas that accretes from a
region of size $l = 2b_{\text{capture}}$ should be
\begin{equation}
L_{\text{turb}} \simeq 1 \times 10^{20}\;\text{cm}^2\;\text{sec}^{-1}
\left(\frac{M}{M_\odot}\right)^{3/4}
\left(\frac{u_\infty}{10\;\text{km\,sec}^{-1}}\right)^{-7/2}. \tag{4.8.10}
\end{equation}

Evidently, the angular momentum may be important in some cases and may not
be in others.
